\section{Discussion}
\label{sec:discussion}

\subsection{Guidance for system integrators}

The measurements translate directly into selection guidance for practitioners
who compose commercial speech services. Above a signal-to-noise ratio of
about 10~dB, every recognition system measured here operates within one
point of its clean error rate on clearly articulated speech, and the choice
among them can be driven by latency, cost, and interface ergonomics rather
than accuracy. The latency data in Table~\ref{tab:latency} then favour the
audio-input LLM systems and the specialised streaming recognisers, which
all answer in well under a second, over the omni-modal path. Between 0 and
5~dB, the product classes separate, and the knee thresholds of
Table~\ref{tab:slopes} express the difference most usefully: the LLM-ASR
systems cross ten percent error 2.2 to 2.5~dB later in Mandarin and 4.4 to
4.8~dB later in English than the specialised streaming system, so in
intermittently noisy environments such as public spaces, vehicles, or
shared offices, the foundation-model recognisers hold the better end of
the accuracy-latency trade-off, with the real-speech caveat that the
advantage is confined to the near-knee regime and reverses for one LLM
system in the hallucination regime (Section~\ref{sec:results:real}). Below 0~dB, every system fails, the LLM
systems by hallucination and the streaming systems by silence, and the
appropriate response is architectural: capture improvement, multi-microphone
processing, or explicit rejection logic. For telephone-band audio, the
surprising result is that no specialist premium exists: the telephone
specialist matches neither the generalists' accuracy in its own band nor
their robustness elsewhere. Its reverberation fragility, 18.2 percent
error against at most 3.2 percent for every other system, proves specific
to the synthetic extreme-DRR response: under three measured room impulse
responses at matched reverberation time it degrades only from 5.1 to 6.3
percent error (Section~\ref{sec:results:real}). Its selection can be
justified only by its 8~kHz interface requirements
and its speed, not by accuracy.

For translation pipelines, the frontend choice dominates the measurements
reported here, and the latency data combine with it into a compact decision
rule. A cascade assembled from the fastest recognition tier and the text
translation engine completes in about 0.7 seconds per sentence on average,
roughly four times as fast as the omni end-to-end path, while the streaming
translation service completes in about half a second and is therefore the
natural choice when latency dominates; when noise robustness dominates, the
same streaming service and the LLM-frontend cascades lead the field, so the
two objectives point to the same shortlist. Swapping the recognition stage
of a cascade from the specialised streaming system to the LLM recogniser is
worth 28 BLEU points at 0~dB (44.1 against 15.8) on an identical translation backend, a far
larger effect than any measurement noise in this study, though the present
design holds the translation engine fixed and therefore cannot rank frontend
changes against engine changes. The pass-through regression of
Section~\ref{sec:results:cac}, at roughly half a BLEU point per percentage
point of recognition error with an intercept that is statistically
indistinguishable from zero and numerically within a point of it, allows an
integrator to rank candidate frontends with recognition-only measurements,
which are cheaper to run and easier to interpret than full-pipeline
evaluations, and to forecast translation quality degradation from a
measured or contracted recognition error rate, with the caveats that the
coefficient is calibrated on one translation engine, that its honest
uncertainty is the condition-level interval rather than the sentence-level
one, and that the linear form saturates in the whole-utterance failure
regime. The end-to-end versus
cascade decision should rest on measured degradation behaviour rather than
architectural intuition: on the point estimate the omni-modal end-to-end
path evaluated here is
simultaneously the slowest and, under severe noise, the most fragile of the
five translation paths, while a dedicated streaming translation service
achieves both the smallest relative loss at 0~dB on synthetic speech and
sub-second latency.

\subsection{Why foundation-model recognisers degrade more slowly}

The robustness advantage of the LLM-ASR systems echoes the finding of
Radford et al.~\cite{radford2023robust} that recognition models trained on
very large and weakly curated corpora acquire substantial noise tolerance
without explicit multi-condition training. Two of our observations sharpen
this picture for API consumers. First, the advantage is strongly
condition-dependent: it is large under energetic babble noise and negligible
under codec, packet-loss, clipping, and tempo distortions, which barely
affect any system. This pattern is consistent with robustness arising from
exposure to diverse recording conditions during pre-training rather than
from any explicit enhancement front-end. Second, the advantage is
language-asymmetric, with Chinese recognition remaining usable at 0~dB
while English recognition of the same systems fails at 37 to 50 percent,
and, as Section~\ref{sec:results:snr} notes, part of the contrast may
reflect the character-level versus word-level scoring scales rather than
raw robustness alone. Providers
disclose neither training mixture nor per-language capacity, so the asymmetry
can only be documented, not explained, at the black-box level; the
documentation itself is actionable, since it implies that multilingual
deployments should be derated per language rather than per provider. It also implies a concrete procurement test: a
noise probe around each system's knee, a few hundred sentences of the target
language in the 0 to 5~dB window, places a candidate service on the
robustness scale documented here, and repeating the probe per deployment
language guards against extrapolating a strong Mandarin result to a weak
English result or the reverse. The direction decomposition of
Section~\ref{sec:results:cascade} extends the same caution to translation
deployments, where the two directions of a bilingual service degraded by
asymmetric amounts, from 1.3 to 1 ratio to 3.2 to 1, under identical noise.

\subsection{Cascade versus end-to-end revisited}

The classical comparison of Bentivogli et al.~\cite{bentivogli2021cascade}
found end-to-end and cascade systems competitive on clean speech with
different error structures. Our measurements extend that comparison to
degraded channels in the black-box setting and qualify it in two ways. The
end-to-end advantage in absolute quality against the same-vendor primary reference is
real but partly a family artefact, since the omni system that translates
directly also shares a generation family with the strongest reference
producer, and the advantage inverts under the secondary reference, as
Section~\ref{sec:results:dual} shows. More importantly, robustness is not an
architectural property but a property of the acoustic frontend wherever it
lives. On synthetic speech the
end-to-end service with the specialised internal frontend is the
most resilient path we measured, while the end-to-end omni path is the most
fragile on the point estimate, and the same frontend ordering that explains cascade behaviour also
explains end-to-end behaviour. The real-speech validation adds a
measurement-side caution to this picture: its translation reference shares
the cascade's engine, so the absolute BLEU gap it reports between the
streaming end-to-end service and the cascade (27 to 28 points on clean
speech) is a same-engine artefact of the kind that
Section~\ref{sec:results:dual} documents, not a translation-quality
ranking. What survives the reference choice is the relative-loss
comparison, in which the end-to-end service loses 41.6 percent of its
clean score under 0~dB real babble against 33.7 percent for the cascade,
a paired contrast of $+7.9$ points whose 95 percent bootstrap interval
($-0.4$ to $+15.7$) still spans zero, so the real-speech data support the
relative-fragility reading as a tendency, without licensing either a
statistically resolved ordering or any absolute claim. For a consumer, the practical
question is not
which architecture to prefer but where the provider placed its best acoustic
model, and the answer is discoverable only by measurement of the kind this
benchmark performs.

\subsection{Methodological implications for black-box evaluation}

Three methodological observations generalise beyond the specific services
tested. The dual-reference protocol of Section~\ref{sec:results:dual}
demonstrates that machine-generated evaluation references can reorder
systems in a way that tracks reference-generator affinity rather than
translation quality, and that degradation-relative conclusions are the
stable currency of black-box comparison. The synthesis-based corpus design
gives black-box studies an exact transcript by construction, which matters
because at sub-percent clean error rates, transcription uncertainty in
recorded corpora would be comparable to the effects being measured. Finally,
the empty-output behaviour documented in Section~\ref{sec:results:main},
where three of the five systems return empty transcripts on a systematic
subset of severely degraded utterances, through two distinct mechanisms,
argues that black-box benchmarks should report response completeness
alongside error rates, since a missing transcript and a wrong transcript
have different downstream failure modes, and a silently empty one differs
operationally from a protocol-level refusal.

\subsection{A worked example of the pass-through rule}

The pass-through relation is most useful when stated as an explicit
calculation, including its intercept and its validity bounds. For the
Qwen-Audio cascade the fitted relation is a loss of $0.44\,e + 0.6$ BLEU
points at recognition error $e$ percent, and for the specialist cascade
$0.54\,e - 1.0$, where the intercepts are point estimates that are
statistically indistinguishable from zero and are carried for arithmetic
consistency with the fitted line. Suppose an integrator operates that
cascade in an environment that
occasionally degrades to 10~dB, at which the recogniser errs at 0.8 percent
on the sweep subset; the predicted translation loss is 0.9 BLEU points,
which is below the clean-condition agreement between the two reference
systems and therefore negligible, and consistent with the measured 1.0
point loss at that condition. If the same deployment used the specialised
streaming recogniser, whose 10~dB error rate is 2.4 percent on the same
sentences, the specialist-calibrated regression predicts $0.54 \times 2.4 -
1.0$, about 0.3 BLEU points, again matching the measured figure within a
point. In the 0~dB regime the rule predicts a loss of $0.44 \times 24.1 +
0.6$, about 11 points, against a measured 9.8, and for the specialist
frontend $0.54 \times 59.2 - 1.0$, about 31 points, against a measured
35.7, so the rule is accurate to within a few points in the collapse
regime but, as Section~\ref{sec:results:cac} quantifies, it is an
empirical interpolation over the measured conditions rather than a law:
its condition-level uncertainty spans 0.23 to 0.45 for the Qwen-Audio
cascade, it is calibrated on one translation engine, and it must not be
extrapolated to recognition errors beyond the measured range, where BLEU
loss saturates while recognition error can grow without bound. The same
arithmetic applied to the measured
error rates of a candidate system at the actual operating signal-to-noise
distribution of a deployment yields an expected translation quality loss
without running the translation stage at all.

\subsection{Implications for evaluation practice and future work}

Beyond the specific findings, the benchmark design carries lessons for
future evaluations of commercial speech intelligence. A benchmark intended
to track black-box services over time should log every raw response rather
than only aggregate scores, because behavioural changes such as the refusal
pattern documented here appear as distributional shifts before they appear
in means. It should treat reference construction as a first-class
experimental variable, since Section~\ref{sec:results:dual} shows that the
reference generator can reorder systems. It should separate degradation
gradients, which generalise across severity, from point estimates, which do
not, and it should report response completeness alongside accuracy. The
present study points to four directions for follow-up work: replication
across providers to separate family effects from provider effects, extension
to spontaneous and accented speech, sweep-based characterisation of the
currently fixed-severity conditions, and human-reference scoring of the
translation stage to anchor the machine-reference comparisons. As foundation
model services begin to expose streaming speech interfaces with embedded
recognition, the measurement approach demonstrated here, treating each
service as an opaque function from degraded audio to text and scoring the
complete pipeline end to end, is, at present, the most direct methodology available for
the comparative robustness questions that practitioners actually face.

\section{Limitations}
\label{sec:limitations}

The benchmark uses synthetic speech from four neural voices, each reading
each sentence once, so voice-level and take-level variability is not
sampled and the corpus is read-style by construction. Synthetic
speech lacks the disfluencies, accents, emotional variation, and channel
noise of spontaneous human recordings, and although the exact-transcript
property is decisive for measurement at sub-percent error floors, clean
error rates measured on synthetic speech should not be read as population
estimates for human speech; the real-speech validation of
Section~\ref{sec:results:real} addresses this concern for read-style
recordings and reproduces the degradation structure, but it covers 150
utterances from two read-speech corpora, not spontaneous or accented
speech, so the clean-condition levels remain indicative rather than
population estimates. The most consequential design
constraint is that all systems, references, and the synthesis come from a
single provider, so product class and vendor are fully confounded: the
class-level regularities reported here are consistent with published
training-scale explanations~\cite{radford2023robust}, but cross-provider
replication is required before they can be attributed to architecture
rather than to one provider's engineering. This revision bounds the
confound with three hold-out measurements --- a locally run open-source
recogniser that replicates the knee-and-collapse and hallucination
structure outside the vendor family, a family-disjoint GLM translation
reference and backend over which the relative-loss ordering is stable, and
a measured-RIR reverb calibration --- but these bound the structural
findings only; the absolute knee locations, slopes, and score levels
remain single-vendor measurements, and the open-source replication used a
small quantised model, chosen for CPU reproducibility rather than
competitive accuracy. The GLM reference is itself an LLM-generated
reference rather than a human one, so absolute translation quality still
lacks a human anchor; human references remain the needed anchor for
absolute translation quality, although the
real-speech validation sidesteps this concern for recognition, where its
references are human transcriptions. The synthetic babble
noise is smoother than real multi-talker cocktail-party interference,
which includes semantic competition and stronger nonstationarity; the
real-babble validation condition confirms that this gap makes synthetic
babble the milder stressor, shifting the degradation knee upward by
roughly 3--6~dB, without changing the system ordering or the collapse
structure. Similarly, the reverberation condition uses a synthetic
exponentially decaying response whose diffuse tail yields a
direct-to-reverberant ratio of $-27.6$~dB, far below the $+7.3$~dB survey
median for everyday spaces; the measured-impulse-response validation
shows real rooms at matched $T_{60}$ to be markedly milder, so the
synthetic reverberation results should be read as a stress case at the
reverberant extreme. Severity levels
for the non-noise degradations were fixed at deployed-but-aggressive values
rather than swept, so only the noise family is characterised by a curve,
and the Opus and narrowband conditions proved to be partially redundant
stressors at the chosen bitrate. Translation quality is measured with BLEU
and chrF against machine-generated references, with the family-affinity
bias that Section~\ref{sec:results:dual} quantifies; no learned
reference-free metric (COMET class) or human judgement was available within
the zero-cost constraint, and human-reference scoring is the single
addition that would most strengthen the absolute claims. The corpus
sentences were authored to be neutral and predictable, so the language-model
prior of the LLM recognisers may partly mask acoustic degradation, and the
numerals-domain gap may partly reflect inverse text normalisation rather
than acoustics, as the normalisation sensitivity analysis suggests.
Finally, commercial APIs are moving targets: the
measurements document the October 2026 versions of the listed services, one
system could not complete its full condition grid before its
free-tier quota was exhausted, and the websocket transport of the streaming
translation service dropped roughly ten percent of its sessions, which is
itself a practically relevant datum
for anyone planning a similar measurement on complimentary quotas.
